\section*{الفصل \ref{ch:b3:planetary-interiors-surfaces} --- البنى الداخلية للكواكب وأسطحها}

\begin{solution}{exo:b3:planetary-interiors-surfaces:1}
$\tfrac52C\bar\rho = 2.5\times0.35\times5429 = 4750$; $x^2 = (4750 -
3300)/(5429 - 3300) = 0.681$، $x = 0.825$، لبٌّ بنصف قطر $0.825\times2440 =
\qty{2010}{km}$; $\rho_c = 3300 + 2129/0.825^3 = \qty{7090}{kg.m^{-3}}$.
\end{solution}

\begin{solution}{exo:b3:planetary-interiors-surfaces:2}
$m = 3500\times\tfrac43\pi(25)^3 = \qty{2.29e8}{kg}$; $E = \tfrac12 m v^2 =
\tfrac12\times\num{2.29e8}\times(\num{2.0e4})^2 = \qty{4.6e16}{J}$، أو
$\num{4.6e16}/\num{4.2e15} = \qty{11}{Mt}$. تعطي القاعدة فوهةً قطرها
$20\times\qty{50}{m} = \qty{1}{km}$.
\end{solution}

\begin{solution}{exo:b3:planetary-interiors-surfaces:3}
$N_1 = 120/\num{100000} = \qty{1.2e-3}{km^{-2}}$، الخطأ النسبي $1/\sqrt{120}
= \qty{9}{\percent}$. على الجزء الخطّي، $T = N_1/c = \num{1.2e-3}/\num{4.875e-4}
= \qty{2.46}{Ga}$ (الحدّ الأُسّي، $\num{8e-7}$، مهمل)،
و$\sigma_T = T/\sqrt{120} = \qty{0.22}{Ga}$.
\end{solution}

\begin{solution}{exo:b3:planetary-interiors-surfaces:4}
المساحة $4\pi(\num{1.8215e6})^2 = \qty{4.17e13}{m^2}$؛ الاستطاعة $2.5\times
\num{4.17e13} = \qty{1.04e14}{W} = \qty{104}{TW}$، أي 2.2 مرّة إجمالي ما تُطلقه الأرض
من سطح أصغر اثنتي عشرة مرّة، و$2.5/0.092 = 27$ مرّة
دفقها.
\end{solution}

\begin{solution}{exo:b3:planetary-interiors-surfaces:5}
(أ)~$\bar\rho = \num{6.417e23}/(\tfrac43\pi(\num{3.3895e6})^3) =
\qty{3934}{kg.m^{-3}}$; $g = GM/R^2 = \qty{3.73}{m.s^{-2}}$; $\rho gR =
\qty{4.97e10}{Pa}$. (ب)~$19\mu/(2\rho gR) = 19\times\num{6.7e10}/\num{9.94e10}
= 12.8$، $k_2 = 1.5/13.8 = 0.109$. (ج)~القيمة المقاسة أكبر بمقدار
2.6 انحراف معياري: يُعطي المريخ قيمةً أكبر من كوكب صلب ذي
صلابة الوشاح، كما كان سيفعل لو احتوى لبًّا سائلًا.
\end{solution}

\begin{solution}{exo:b3:planetary-interiors-surfaces:6}
(أ)~القمر: $4\times9.8/1.62 = \qty{24}{km}$ في الصخر البلّوري،
\qty{12}{km} في الرواسب؛ المريخ وعطارد: $4\times9.8/3.7 =
\qty{10.6}{km}$ و\qty{5.3}{km}. تقع قيمة القمر البالغة \qty{20}{km} بين
قيمتَي توقّعها: القاعدة صحيحة إلى بضعة عشرات في المئة. (ب)~$Y =
\rho gD/3 = 2600\times9.8\times4000/3 = \qty{34}{MPa}$، مقاومةَ
الصخر المتشقّق لا البلّورة السليمة.
\end{solution}

\begin{solution}{exo:b3:planetary-interiors-surfaces:7}
المريخ: $\num{3.22e-14}(e^{6.93T} - 1) + \num{4.875e-4}\,T = \num{2.0e-3}$؛
التخمين الخطّي $4.1$~Ga أقدم بكثير من اللازم، إذ الحدّ الأُسّي كبير
هناك؛ تعطي طريقة التنصيف $T = \qty{3.34}{Ga}$ (الحدّ الأُسّي
$\num{3.6e-4}$، الحدّ الخطّي $\num{1.63e-3}$). تعطي الدالّة القمرية
\qty{2.39}{Ga}. كلّ دالّة مُعايَرة لجرم واحد: معدّل
الاصطدامات وحجم الفوهات التي يُحدثها مقذوف معيّن تختلف
بين القمر والمريخ.
\end{solution}

\begin{solution}{exo:b3:planetary-interiors-surfaces:8}
(أ)~$3.551181/1.769138 = 2.0073$ و$7.154553/3.551181 = 2.0147$، قريبة من
$2$. (ب)~$0.565247 - 0.844789 + 0.279542 = 0.000000$ في اليوم، بدقّة
$\num{1e-7}$: الحركات المتوسّطة الثلاث مقفلة، وهو رنين يحفظ
المواضع النسبية للأقمار عند كلّ اقتران. (ج)~تتكرّر قوى الجذب عند
النقاط نفسها من المدار، فتتراكم بدل أن تتلاشى؛
فتغذّي اللامركزية بالسرعة نفسها التي يخمدها بها المدّ.
\end{solution}

\begin{solution}{exo:b3:planetary-interiors-surfaces:9}
$n = 2\pi/(1.769138\times86400) = \qty{4.111e-5}{s^{-1}}$, $nR =
\qty{74.9}{m.s^{-1}}$, $(nR)^5/G = \qty{3.53e19}{W}$;
$\tfrac{21}{2}e^2(nR)^5/G = 10.5\times\num{1.6e-5}\times\num{3.53e19} =
\qty{5.9e15}{W}$ لكلّ وحدة من $k_2/Q$. ثمّ $k_2/Q =
\num{1.04e14}/\num{5.9e15} = 0.018$.
\end{solution}

\begin{solution}{exo:b3:planetary-interiors-surfaces:10}
(أ)~$x^2 = (\tfrac52C\bar\rho - \rho_m)/(\bar\rho - \rho_m)$ هي خطّية في
$C$، ميلها $\tfrac52\bar\rho/(\bar\rho - \rho_m)$. (ب)~الميل يساوي
$2.5\times3340/52 = 161$: يتغيّر $x^2$ بمقدار $0.161$ مقابل $0.001$ من $C$.
من أجل $C = 0.395$، $x^2 = 0.036 + 0.161 = 0.197$، $x = 0.44$، $\rho_c =
\qty{3880}{kg.m^{-3}}$؛ ومن أجل $C = 0.393$، $x^2 < 0$: لا لبّ أصلًا.
(ج)~قرب $C = \tfrac25$ عزم القصور الذاتي لا يكاد يتأثّر بلبٍّ صغير:
رقمه المعنوي الثالث يحسم بين لبٍّ كبير خفيف وعدم وجود لبّ.
تلزم الموجات الزلزالية والمدّ والمجال المغناطيسي.
\end{solution}

\begin{solution}{exo:b3:planetary-interiors-surfaces:11}
(أ)~$\sigma_T = N_1/(10\,N_1')$: عند \qty{1.0}{Ga}، $N_1 \approx N_1' \approx
\num{4.875e-4}$، $\sigma_T = \qty{0.10}{Ga}$؛ وعند \qty{3.8}{Ga}، $N_1 =
\num{1.07e-2}$ و$N_1' = \num{6.15e-2}$، $\sigma_T = \qty{0.017}{Ga}$،
أصغر بستّ مرّات. (ب)~مع $N_1 = cT$، $N_1' = c$ و$\sigma_T =
cT/(\sqrt n\,c) = T/\sqrt n$. (ج)~دالّة الكرونولوجيا نفسها مشوبة
بعدم اليقين: إذ تستند إلى عدد قليل من المواقع القمرية المؤرَّخة وإلى التدرّج من القمر
إلى المريخ، وهي أخطاء لا يُخفّفها أيّ عدّ.
\end{solution}

\begin{solution}{exo:b3:planetary-interiors-surfaces:12}
(أ)~النسبة هي $(n_ER_E/n_IR_I)^5(e_E/e_I)^2 = (0.4982\times0.8569)^5
\times(0.009/0.004)^2 = 0.0142\times5.06 = 0.072$: نحو
\qty{7.5e12}{W}، أو $\num{7.5e12}/(4\pi(\num{1.5608e6})^2) =
\qty{0.24}{W.m^{-2}}$. (ب)~يفصل \omterm{def:b3:planetary-interiors-surfaces:ocean-world}{المحيط تحت السطحي} القشرةَ الجليدية عن
الباطن: تنثني القشرة أكثر بكثير مما ينثني قمرٌ صلب، ويتبدّد جليدها
بطريقة تختلف عن صخر آيو الساخن، فيتغيّر كلٌّ من $k_2$ و$Q$.
\end{solution}

\begin{solution}{pb:b3:planetary-interiors-surfaces:1}
\textbf{1.} $\bar\rho = \num{6.4169e23}/(\tfrac43\pi(\num{3.3895e6})^3) =
\qty{3934}{kg.m^{-3}}$.

\textbf{2.} $\tfrac52C\bar\rho = 3600$; $x^2 = (3600 - 3500)/(3934 - 3500) =
0.230$، $x = 0.479$، لبٌّ بنصف قطر \qty{1620}{km}؛ $\rho_c = 3500 +
434/0.479^3 = \qty{7450}{kg.m^{-3}}$.

\textbf{3.} $x = 1830/3389.5 = 0.540$؛ $\rho_m = (3600 - 0.2915\times
3934)/(1 - 0.2915) = \qty{3460}{kg.m^{-3}}$، $\rho_c = 3462 +
472/0.1574 = \qty{6460}{kg.m^{-3}}$.

\textbf{4.} وهو أدنى بالثلث. يكون الضغط في لبّ المرّيخ
أدنى بكثير منه في لبّ الأرض، فيكون معدنه أقلّ انضغاطًا؛ ولبٌّ بهذه الخفّة يوحي
بعناصر أخفّ بكثير مختلطة بالحديد.

\textbf{5.} جاء \omterm{def:b3:planetary-interiors-surfaces:love-number}{عدد لوف} (2003) قبل مهمّة إنسايت (2021) بسنوات طويلة:
أظهر بالفعل أنّ اللبّ الخارجي على الأقلّ سائل، وحدّد نصف قطره
بين 1520 و\qty{1840}{km}، وهو مدًى يشمل
القيمة الزلزالية.

\textbf{6.} $1 - (1 - 20/3389.5)^3 = \qty{1.8}{\percent}$ إلى $1 - (1 -
37/3389.5)^3 = \qty{3.2}{\percent}$ من الحجم، وأقلّ من الكتلة:
فهو أخفّ من الوشاح، فيخفّض قليلًا كثافة الوشاح الفعّالة،
وهو أثر أصغر من عدم اليقين في النموذج.

\textbf{7.} من \qty{1}{km}: 54؛ من 1.41: 16؛ من 2: 10؛ من 2.83: 5؛
من 4، 5.66، 8 و11.3: 2؛ من 16: 1.

\textbf{8.} $N = 54/80000 = \qty{6.75e-4}{km^{-2}}$، $\sigma = \sqrt{54}/80000
= \qty{0.92e-4}{km^{-2}}$ (\qty{14}{\percent}).

\textbf{9.} المتوقَّع $54\times0.355 = 19$ فوهة من \qty{1.41}{km}
و$54\times0.147 = 7.9$ من \qty{2}{km}؛ والمعدود 16 و10،
ضمن انحراف معياري واحد ($\sqrt{19} = 4.4$، $\sqrt{7.9} = 2.8$): متّسق.

\textbf{10.} لا: $0.072/\num{6.75e-4} = 107$ مرّة أدنى منه. العدد
يسجّل الزمن.

\textbf{11.} الفئات متساوية العرض في $\log D$ تغطّي المدى بانتظام وتُبقي
أعداد الفوهات متقاربة في فئات التوزيع الشديد الانحدار.

\textbf{12.} هي فوهات ثانوية، تكوّنت من كتل قُذفت من فوهة أولية كبيرة:
ليست اصطدامات مستقلّة قادمة من الفضاء، ولو كانت كذلك لأضافت
فوهات لا تحسبها \omterm{def:b3:planetary-interiors-surfaces:crater-dating}{دالة كرونولوجيا الفوهات}.

\textbf{13.} $T = \num{6.75e-4}/\num{4.875e-4} = \qty{1.385}{Ga}$؛ والحدّ
الأسّي، $\num{3.22e-14}(e^{9.6} - 1) = \num{4.7e-10}$، هو جزء من
مليون من $N_1$.

\textbf{14.} $\sigma_T = T/\sqrt{54} = \qty{0.19}{Ga}$: $1.38 \pm 0.19$~Ga.

\textbf{15.} $\num{6.75e-4}/\num{8.38e-4} = \qty{0.81}{Ga}$، أصغر بنحو
النصف تقريبًا.

\textbf{16.} $31 + 12 + 7 + 2 + 0 + 2 + 2 = 56$ فوهة؛ $N(\qty{2}{km}) =
56/120000 = \qty{4.67e-4}{km^{-2}}$؛ $N_1 = \num{4.67e-4}/0.1473 =
\qty{3.17e-3}{km^{-2}}$. بطريقة التنصيف الثنائي على $\num{3.22e-14}(e^{6.93T} - 1) +
\num{4.875e-4}\,T = \num{3.17e-3}$ تُعطي $T = \qty{3.54}{Ga}$ (الحدّ الأسّي
\num{1.45e-3}، والحدّ الخطّي \num{1.73e-3}).

\textbf{17.} $N_1' = \num{3.22e-14}\times6.93\,e^{6.93\times3.54} +
\num{4.875e-4} = \num{1.05e-2}$؛ $\sigma_T = \num{3.17e-3}/(\sqrt{56}\times
\num{1.05e-2}) = \qty{0.040}{Ga}$.

\textbf{18.} الحالة A البسيطة: \qty{14}{\percent}؛ والحالة B البسيطة: \qty{1.1}{\percent}. في
الجزء المبكّر الشديد الانحدار من دالة الكرونولوجيا، يعني تغيّرًا كبيرًا في كثافة الفوهات
تغيّرًا صغيرًا في العمر؛ وعلى الجزء الخطّي يكون العمر غير مؤكّد بقدر
العدّ نفسه.

\textbf{19.} $120000\times\num{3.17e-3}\times(1 - 0.1473) = 324$؛ فهو
يحمل $48 + 8 = 56$.

\textbf{20.} $N_1 = 56/(120000\times0.8527) = \qty{5.47e-4}{km^{-2}}$،
$T_r = \num{5.47e-4}/\num{4.875e-4} = \qty{1.12}{Ga}$.

\textbf{21.} $\sigma = 1.12/\sqrt{56} = \qty{0.15}{Ga}$.

\textbf{22.} فوهة قطرها \qty{2}{km} عمقها \qty{400}{m}: كانت الحمم على
الأقلّ بهذا السُّمك. والفوهات الأكبر أعمق، وترتفع حوافّها أعلى؛
امتلأت جزئيًّا وما زالت تظهر.

\textbf{23.} $1.38 \pm 0.19$ و$1.12 \pm 0.15$~Ga يختلفان بمقدار \qty{0.26}{Ga}،
بنحو انحراف معياري مشترك واحد ($\sqrt{0.19^2 + 0.15^2} = 0.24$):
تسمح الأعداد بحدثٍ واحد، وإن كانت لا تُثبته.

\textbf{24.} يبلغ العمر النموذجي للسهل الحمَمي الفتيّ A
\textbf{1.4~مليار~سنة} ($1.38 \pm 0.19$~Ga).
\end{solution}
